\documentclass[11pt]{article}

% Keep byte-for-byte output stable across rebuild directories and wall times.
\pdfinfoomitdate=1
\pdftrailerid{}
\pdfsuppressptexinfo=15

\usepackage[T1]{fontenc}
\usepackage{amsmath,amssymb,amsthm,mathtools}
\usepackage{enumitem}
\usepackage{booktabs}
\usepackage{array,tabularx,longtable}
\usepackage{seqsplit}
\usepackage[hidelinks]{hyperref}
\hypersetup{
  pdftitle={Real Closedness of Hahn and Puiseux Series Fields},
  pdfauthor={selpo}
}

\newcommand{\Hahn}{k((\Gamma))}
\newcommand{\Puiseux}{\mathcal P(k)}
\newcommand{\Val}{v}
\newcommand{\Supp}{\operatorname{supp}}
\newcommand{\lc}{\operatorname{lc}}
\newcommand{\Init}{\operatorname{in}}
\newcommand{\Res}{\operatorname{res}}
\newcommand{\Ocal}{\mathcal O}
\newcommand{\Odd}{\operatorname{Odd}}
\newcommand{\lean}[1]{\texttt{\expandafter\seqsplit\expandafter{\detokenize{#1}}}}
\newcommand{\namedref}[2]{\hyperref[#2]{#1~\ref*{#2}}}
\newcommand{\claimref}[3]{%
  \hyperref[#1]{\textbf{#2~\ref*{#1}:} #3}}

\newtheoremstyle{proofnote}
  {6pt}{6pt}{\normalfont}{}{\bfseries}{.}{.5em}
  {\thmname{#1}\thmnumber{ #2}\thmnote{ (#3)}}
\theoremstyle{proofnote}
\newtheorem{theorem}{Theorem}[section]
\newtheorem{proposition}[theorem]{Proposition}
\newtheorem{lemma}[theorem]{Lemma}
\newtheorem{corollary}[theorem]{Corollary}
\newtheorem{definition}[theorem]{Definition}

\title{Real Closedness of Hahn and Puiseux Series Fields}
\author{selpo}
\date{}

\begin{document}
\maketitle

\begin{abstract}
  Let \(k\) be a real closed ordered field and let \(\Gamma\) be a divisible
  ordered abelian group.  We prove that the Hahn field \(k((\Gamma))\) is
  real closed, and that its bounded-denominator subfield with
  \(\Gamma=\mathbb Q\) is real closed as well.  The first proof combines
  valuation-theoretic normalization with transfinite Newton corrections.  At
  a limit stage the corrections are summed over the whole ordinal initial
  segment, and a set-theoretic size argument rules out an indefinitely
  increasing sequence of exponents of the added monomials.  The second proof
  keeps the corrections inside one rational
  denominator lattice.  A correction whose selected root retains the current
  multiplicity cannot leave that lattice,
  so an infinite sequence of such exponents is unbounded and its Hahn sum is
  a Puiseux root.
\end{abstract}

\tableofcontents

\section*{Proof overview}
\addcontentsline{toc}{section}{Proof overview}

The proof is organized as follows.  The ordered-field criterion reduces
real closedness to square roots of nonnegative elements and roots of
odd-degree polynomials.  For Hahn series, the odd-degree problem is converted
to coefficient conditions over the valuation ring.  Multiplicity one is
solved by ordinal recursion.  At larger odd multiplicity, a Newton correction
either reduces the selected root multiplicity or preserves it; the latter
case is continued over complete ordinal initial segments; a sufficiently long
continuation would give an impossible injection into \(\Gamma\).  For
Puiseux series, the root construction is repeated while preserving one common
denominator lattice, using Newton iteration for a simple root over the
corresponding power-series ring.

\section{Main statements and the ordered-field criterion}
\label{sec:main}

An ordered field is called \emph{real closed} if it has no proper algebraic
extension to which its order extends; equivalently, adjoining a square root
of \(-1\) gives an algebraically closed field.  Let \(k\) be a real closed
ordered field and let \(\Gamma\) be a divisible linearly ordered additive
group.  Here divisible means that, for every
\(\gamma\in\Gamma\) and every integer \(n\ge1\), there is a (necessarily
unique) \(\delta\in\Gamma\) with \(n\delta=\gamma\).  A Hahn series over
\(k\) with exponents in \(\Gamma\) is a formal sum
\[
  x=\sum_{\gamma\in\Gamma} a_\gamma t^\gamma
\]
whose support
\[
  \Supp(x)=\{\gamma\in\Gamma:a_\gamma\ne0\}
\]
is well ordered, meaning that every nonempty subset of the support has a
least element.  Addition is coefficientwise and multiplication is the
convolution product.  The well-ordering condition is exactly what makes
every coefficient of a product a finite sum and keeps the support
well ordered.  The resulting field is denoted by
\[
  \Hahn=k((\Gamma)).
\]
It is ordered lexicographically: the sign of a nonzero series is the sign of
the coefficient at the least exponent in its support.

\begin{theorem}[Hahn--Kaplansky]
\label{thm:hahn-real-closed}
  The ordered field \(k((\Gamma))\) is real closed.
\end{theorem}

For \(\Gamma=\mathbb Q\), define the bounded-denominator Puiseux field by
\[
  \Puiseux
  =
  \left\{
    x\in k((\mathbb Q)):
    \text{there is }N\ge1\text{ with }
    \Supp(x)\subseteq \frac1N\mathbb Z
  \right\}.
\]
The sum, product, and inverse of two elements with a common denominator
bound have the same property, so \(\Puiseux\) is a subfield of
\(k((\mathbb Q))\).

\begin{theorem}[Puiseux real closedness]
\label{thm:puiseux-real-closed}
  The ordered field \(\Puiseux\) is real closed.
\end{theorem}

We use the standard characterization of real closed ordered fields; the
valuation-theoretic background is standard, see \cite{Kaplansky,EnglerPrestel}.
\begin{proposition}[Ordered-field criterion]
\label{prop:real-closed-criterion}
  An ordered field \(K\) is real closed if the following two conditions hold:
  \begin{enumerate}[label=(\roman*)]
    \item every \(x\ge0\) in \(K\) is a square;
    \item every polynomial in \(K[X]\) of odd degree has a root in \(K\).
  \end{enumerate}
\end{proposition}

Thus both theorems reduce to the same two tasks.  The square-root task is
local at the least exponent of a series.  The odd-degree root task is
global: Newton corrections have to be summed, and their exponents have to
be shown to remain under control.

\section{Valuation preliminaries and square roots}
\label{sec:square-roots}

For \(x\ne0\) in \(k((\Gamma))\), put
\[
  \Val(x)=\min\Supp(x),\qquad
  \lc(x)=a_{\Val(x)},
\]
and put \(\Val(0)=+\infty\).  The least-exponent rule gives
\[
  \Val(xy)=\Val(x)+\Val(y),\qquad
  \Val(x+y)\ge
  \min\{\Val(x),\Val(y)\},
\]
with equality in the second formula when the two valuations are different.
The valuation ring and its maximal ideal are
\[
  \Ocal=\{x:\Val(x)\ge0\},
  \qquad
  \mathfrak m=\{x:\Val(x)>0\}.
\]
The residue of an element of \(\Ocal\) is its coefficient at exponent zero.
An element of \(\Ocal\) is a \emph{unit} when its inverse also belongs to
\(\Ocal\); equivalently, it has valuation zero.  We call a series, or a
polynomial coefficient, \emph{integral} when it belongs to \(\Ocal\).

A family \((A_i)_{i\in I}\) of Hahn series is \emph{Hahn summable} if
\(\bigcup_i\Supp(A_i)\) is well ordered and, for every \(\gamma\in\Gamma\),
only finitely many \(i\) have \([t^\gamma]A_i\ne0\).  Its Hahn sum is then
defined coefficientwise by those finite sums.

\begin{lemma}[Leading-term normalization]
\label{lem:leading-term}
  If \(x>0\) in \(k((\Gamma))\), then
  \[
    x=c\,t^\gamma(1+w)
  \]
  for some \(\gamma\in\Gamma\), \(c\in k\) with \(c>0\), and
  \(w\in\mathfrak m\).
\end{lemma}
\begin{proof}
  Take \(\gamma=\Val(x)\) and \(c=\lc(x)\).  The lexicographic order gives
  \(c>0\).  Put \(w=c^{-1}t^{-\gamma}x-1\).  The series
  \(c^{-1}t^{-\gamma}x\) has valuation zero and leading coefficient one;
  hence its difference from \(1\) has positive valuation.  Thus
  \(w\in\mathfrak m\), and the stated identity follows from the definition
  of \(w\).
\end{proof}

An element \(u\) with \(u-1\in\mathfrak m\), equivalently an element of the
form \(1+w\) with \(w\in\mathfrak m\), is called a \emph{\(1\)-unit}.

We use Neumann's support lemma: if \(S\) is a well-ordered subset of the
positive cone of an ordered abelian group, then the additive monoid generated
by \(S\) is well ordered, and a fixed group element has only finitely many
representations as a sum of a finite sequence from \(S\).  Consequently the
family of powers \((w^n)_{n\ge0}\) is Hahn summable whenever
\(\Supp(w)\subseteq\Gamma_{>0}\).  This is the standard support lemma used in
the construction of the Hahn field; see \cite{EnglerPrestel}.

\begin{lemma}[A \(1\)-unit has a binomial square root]
\label{lem:positive-unit-square}
  If \(u\in k((\Gamma))\) satisfies \(\Val(u-1)>0\), then there is
  \(s\in k((\Gamma))\) with \(u=s^2\).
\end{lemma}
\begin{proof}
  Put \(w=u-1\).  If \(w=0\), take \(s=1\).  Otherwise let
  \(\epsilon=\min\Supp(w)>0\).  The binomial series
  \[
    s=\sum_{n\ge0}\binom{1/2}{n}w^n
  \]
  is a Hahn series.  Indeed, Neumann's support lemma applies to the
  well-ordered positive set \(\Supp(w)\), so the union of the supports of the
  powers is well ordered and only finitely many powers contribute to any
  fixed exponent.  Consequently the
  binomial identity can be checked coefficient by coefficient.  It gives
  \(s^2=1+w=u\).
\end{proof}

\begin{proposition}[Nonnegative elements are squares]
\label{prop:hahn-squares}
  Every \(x\ge0\) in \(k((\Gamma))\) is a square.
\end{proposition}
\begin{proof}
  The assertion is clear for \(x=0\).  For \(x>0\), write
  \(x=c\,t^\gamma(1+w)\) as in \namedref{Lemma}{lem:leading-term}.
  Divisibility of \(\Gamma\) gives \(\delta\) with
  \(2\delta=\gamma\), and real closedness of \(k\) gives
  \(d\in k^\times\) with \(d^2=c\).  By
  \namedref{Lemma}{lem:positive-unit-square}, choose \(s^2=1+w\).  Then
  \[
    x=(d\,t^\delta s)^2.
  \]
\end{proof}

\section{Newton reduction for odd-degree polynomials}
\label{sec:newton-reduction}

Let \(F(X)=\sum_i a_iX^i\in k((\Gamma))[X]\) and let
\(\delta\in\Gamma\).  In the root construction \(\delta>0\) will be the
exponent of a monomial correction \(ct^\delta\).  The Newton weight of the
\(i\)-th term is
\[
  W_\delta(i)=\Val(a_i)+i\delta.
\]
For \(\mu\in\Gamma\), write \([t^\eta]a\) for the coefficient of a Hahn
series \(a\) at exponent \(\eta\), and define the Newton polynomial on the
weight level \(\mu\) by
\[
  \Init_{\delta,\mu}(F)(T)
  =
  \sum_{\substack{i:\ a_i\ne0\\ W_\delta(i)=\mu}}
    [t^{\mu-i\delta}]a_i\,T^i
  \in k[T].
\]
Its nonzero terms are exactly the coefficient points on the line of weight
\(\mu\).  If \(F\ne0\) and
\(\mu_\delta=\min_{a_i\ne0}W_\delta(i)\), then
\(\Init_{\delta,\mu_\delta}(F)\) is the usual initial polynomial on the
lowest Newton line; its coefficient in degree \(i\) is \(\lc(a_i)\) whenever
\(W_\delta(i)=\mu_\delta\).

For a nonzero polynomial \(P\in k[T]\), the \emph{multiplicity} of a root
\(c\) is the largest \(r\in\mathbb N\) for which \((T-c)^r\) divides
\(P\).  Equivalently, \(P(c)=P'(c)=\cdots=P^{(r-1)}(c)=0\) and
\(P^{(r)}(c)\ne0\).

\begin{lemma}[Normalization at an odd-multiplicity Newton root]
\label{lem:newton-correction}
  Let \(\mu\in\Gamma\), assume \(\mu\le W_\delta(i)\) for every \(i\), and
  suppose \(c\in k\) is a root of \(\Init_{\delta,\mu}(F)\) of positive odd
  multiplicity \(r\).  Put
  \[
    G(Z)=t^{-\mu}F\bigl(t^\delta(Z+c)\bigr)=\sum_j g_jZ^j.
  \]
  Then
  \[
    \Val(g_j)\ge0,\qquad
    \Val(g_j)>0\ (j<r),\qquad
    \Val(g_r)=0.
  \]
  Thus \(G\in\Ocal[Z]\), its degree-\(r\) coefficient is a unit, and all
  coefficients below degree \(r\) lie in \(\mathfrak m\).
\end{lemma}
\begin{proof}
  Write \(F(c\,t^\delta+Y)=\sum_j b_jY^j\).  Expanding gives
  \[
    F(c\,t^\delta+Y)
    =\sum_i a_i\sum_{j=0}^i\binom ij
       c^{\,i-j}t^{(i-j)\delta}Y^j.
  \]
  The contribution to \(b_j\) at exponent \(\mu-j\delta\) is
  \(t^{\mu-j\delta}\Init_{\delta,\mu}(F)^{(j)}(c)/j!\); after assigning
  weight \(\delta\) to \(Y\), the corresponding term has total weight
  \(\mu\).  The multiplicity-\(r\)
  assumption makes these coefficients vanish for \(j<r\) and makes the
  coefficient in degree \(r\) nonzero.  All other contributions have
  strictly larger weighted value.  Therefore
  \[
    \Val(b_j)+j\delta\ge\mu,\qquad
    \Val(b_j)+j\delta>\mu\ (j<r),\qquad
    \Val(b_r)+r\delta=\mu.
  \]

  Since \(g_j=t^{j\delta-\mu}b_j\), the three weighted inequalities
  are exactly the three displayed valuation statements for \(G\).  Together
  with the assumed positivity and oddness of \(r\), these are precisely the
  asserted coefficient conditions.
\end{proof}

\begin{definition}[Normalized odd Newton problem]
\label{def:odd-newton-problem}
  A normalized odd Newton problem of multiplicity \(m\) consists of a
  polynomial \(F\in k((\Gamma))[X]\), with \(m>0\) odd, all of whose
  coefficients are integral, whose coefficient of \(X^m\) is a unit, and
  whose coefficients of degree below \(m\) lie in \(\mathfrak m\).
  Equivalently, the coefficientwise lift of \(F\) to \(\Ocal[X]\) has these
  last two properties; we use that lift without changing notation whenever
  the valuation-ring formulation is needed.
\end{definition}

Here a \emph{translation} is a substitution \(X\mapsto X+a\), and a
\emph{monomial dilation} is a substitution \(X\mapsto t^\delta X\), possibly
followed by multiplication of the polynomial by a nonzero monomial
\(t^\eta\).  A solution of a normalized problem is an element
\(y\in\mathfrak m\) with \(F(y)=0\).  If the problem was obtained by finitely
many such transformations, their inverses carry \(y\) to a root of the
original polynomial.

The next lemma is the point at which real closedness of the residue field
enters the Newton argument.

\begin{lemma}[An odd initial polynomial has a nonzero odd-multiplicity root]
\label{lem:proper-odd-root}
  Let \(I\in k[T]\) have odd degree \(n\) and nonzero constant term.
  Then \(I\) has a nonzero root of odd multiplicity \(r\), with
  \(1\le r\le n\).
\end{lemma}
\begin{proof}
  Over a real closed field every irreducible factor has degree one or two.
  Since \(\deg I\) is odd, some linear factor occurs with odd multiplicity;
  let its root be \(a\) and its multiplicity be \(r\).  The nonzero
  constant term says that \(0\) is not a root, hence \(a\ne0\).
  Positivity of root multiplicity and the elementary bound
  \(r\le\deg I=n\) give the remaining assertions.
\end{proof}

For a normalized problem of multiplicity \(m\) with
\(\Val(a_0)=m\delta\), we call it the \emph{lower-edge case at
\(\delta\)} when some \(j<m\) satisfies
\(a_j\ne0\) and \(\Val(a_j)+j\delta<m\delta\).  Thus a coefficient point
lies below the line through the constant and degree-\(m\) points.  A
\emph{proper multiplicity} means a multiplicity strictly smaller than
\(m\).

\begin{proposition}[Lower Newton edges are resolved by multiplicity descent]
\label{prop:lower-edge}
  Write \(F(X)=\sum_i a_iX^i\in k((\Gamma))[X]\).  Let \(m\) be odd (hence
  positive) and assume
  \[
    \Val(a_m)=0,\qquad
    \Val(a_i)>0\quad(0\le i<m),\qquad
    \Val(a_i)\ge0\quad(i>m).
  \]
  Suppose every normalized odd Newton problem of odd multiplicity \(r<m\)
  has a solution.  Assume that for some \(\delta\in\Gamma\)
  \[
    \Val(a_0)=m\delta
  \]
  and that there is an index \(j<m\) with \(a_j\ne0\) and
  \[
    \Val(a_j)+j\delta<m\delta.
    \tag{LE}\label{eq:lower-edge-below}
  \]
  Then \(F\) has a root of positive valuation.
\end{proposition}
\begin{proof}
  For a slope \(s\), call
  \(w_s(i)=\Val(a_i)+is\) the Newton weight of the \(i\)-th coefficient.
  Thus \eqref{eq:lower-edge-below} says exactly that the point belonging to
  \(a_jX^j\) lies strictly below the line of slope \(-\delta\) through the
  constant and degree-\(m\) points: both endpoints have weight \(m\delta\).

  The lower point also proves that \(\delta>0\): rearranging its strict
  inequality gives
  \(0<\Val(a_j)<(m-j)\delta\), and \(m-j>0\).

  Define the first lower slope by
  \[
    \theta=
    \min_{\substack{0\le i<m\\a_i\ne0}}
       \frac{\Val(a_i)}{m-i}.
  \]
  The set is nonempty by \eqref{eq:lower-edge-below}.  Every number in it is
  positive because the coefficients below degree \(m\) lie in
  \(\mathfrak m\), while the index \(j\) in \eqref{eq:lower-edge-below}
  gives \(\theta<\delta\).  Hence
  \[
    0<\theta<\delta.
    \tag{S}\label{eq:lower-edge-theta}
  \]

  Consider
  \[
    \widetilde F_\theta(T)=t^{-m\theta}F(t^\theta T)
      =\sum_i t^{(i-m)\theta}a_iT^i.
  \]
  If \(i<m\), the definition of \(\theta\) gives
  \(\Val(a_i)+(i-m)\theta\ge0\).  The same inequality is immediate for
  \(i=m\), and for \(i>m\) it follows from \(a_i\in\Ocal\) and
  \(\theta>0\).  Thus \(\widetilde F_\theta\in\Ocal[T]\), and we may reduce
  it modulo \(\mathfrak m\):
  \[
    Q(T)=\overline{\widetilde F_\theta(T)}\in k[T].
  \]
  The coefficient of \(T^m\) in \(Q\) is nonzero because
  \(\Val(a_m)=0\).  All coefficients above degree \(m\) vanish after
  reduction because their valuations are strictly positive.  Hence
  \(\deg Q=m\).  An index attaining the minimum defining \(\theta\) gives a
  nonzero coefficient of \(Q\) below degree \(m\).  Finally,
  \[
    \Val(a_0)-m\theta=m(\delta-\theta)>0
  \]
  by \eqref{eq:lower-edge-theta}, so \(Q(0)=0\).

  By the same linear--quadratic factorization argument as in
  \namedref{Lemma}{lem:proper-odd-root}, the odd-degree polynomial \(Q\)
  has a root \(a\) of positive odd multiplicity \(r\); here we do not need
  that the root be nonzero.  Necessarily \(r<m\).  Indeed, if \(r=m\), then
  degree comparison gives \(Q(T)=u(T-a)^m\) for some \(u\ne0\).  Since
  \(Q(0)=0\), this forces \(a=0\), so \(Q=uT^m\), contradicting the
  nonzero coefficient of degree \(j<m\).

  Now translate the root \(a\) to zero:
  \[
    H(Z)=t^{-m\theta}F\bigl(t^\theta(Z+a)\bigr).
  \]
  By \namedref{Lemma}{lem:newton-correction}, \(H\) is a normalized odd
  Newton problem of multiplicity \(r<m\).  The induction hypothesis gives
  \(z\in\mathfrak m\) with \(H(z)=0\).  Put
  \[
    y=t^\theta(z+a).
  \]
  Then \(F(y)=t^{m\theta}H(z)=0\).  Moreover \(z\in\mathfrak m\) and
  \(a\in k\) imply \(\Val(z+a)\ge0\); together with \(\theta>0\), this gives
  \(\Val(y)=\theta+\Val(z+a)>0\).  Thus \(y\) is the required root.
\end{proof}

For \(m>1\), the reduction of a normalized problem modulo \(\mathfrak m\)
has zero as a root of multiplicity exactly \(m\): all coefficients below
degree \(m\) vanish, while the degree-\(m\) coefficient does not.  We call
this the \emph{same-multiplicity zero branch}.  The phrase records this
specific residue root and its multiplicity; it does not assert that the
original polynomial already has a root.

\section{Transfinite Newton continuation}
\label{sec:transfinite}

We now solve the remaining same-multiplicity zero branch by a transfinite
sequence of Newton corrections.

\begin{definition}[Ordinal terminology]
\label{def:ordinal-terminology}
  A well-order is a linear order in which every nonempty subset has a least
  element.  An ordinal is the order type of a well-ordered set.  If
  \(\lambda\) is an ordinal, its initial segment before \(\alpha<\lambda\)
  consists of all \(\beta<\alpha\).  A successor ordinal has the form
  \(\alpha+1\); a nonzero ordinal with no greatest element is a limit
  ordinal.  Transfinite recursion defines an object at a successor stage from
  the preceding object and at a limit stage from the whole earlier initial
  segment.  We write \(\mathrm{Ord}\) for the class of all ordinals.
\end{definition}

At a Newton stage \(\alpha\), we retain an approximation
\(x_\alpha\in\mathfrak m\) and write
\[
  F_\alpha(Z)=F(x_\alpha+Z)=\sum_i b_iZ^i.
\]
We say that this stage still has multiplicity \(m\) when \(F_\alpha\) is a
normalized odd Newton problem of multiplicity \(m\) in the sense of
\namedref{Definition}{def:odd-newton-problem}.  This is only terminology for
the displayed translated polynomial, not an additional kind of Newton
process.

\begin{lemma}[Successor correction]
\label{lem:successor-correction}
  Let \(F\in\Ocal[X]\), \(m\in\mathbb N\), and \(x\in\mathfrak m\).  Write
  \(F_x(Z)=F(x+Z)=\sum_i b_iZ^i\), and suppose
  \(\Val(F(x))=\gamma>0\), \(\rho>0\), and \(m\rho=\gamma\).  Assume every
  Newton weight of \(F_x\) at \(\rho\) is at least \(\gamma\), and let
  \(c\in k^\times\) be a root of \(\Init_{\rho,\gamma}(F_x)\).  Then
  \[
    x'=x+c t^\rho\in\mathfrak m,
    \qquad \Val(x'-x)=\rho,
  \]
  and \(\Val(F(x'))>\gamma=\Val(F(x))\).
\end{lemma}
\begin{proof}
  Both \(x\) and \(ct^\rho\) have positive valuation, so
  \(x'\in\mathfrak m\), while \(c\ne0\) gives
  \(\Val(x'-x)=\rho\).  In the expansion of \(F_x(ct^\rho)\), all terms
  have weight at least \(\gamma\), and the coefficient at weight
  \(\gamma\) is precisely \(\Init_{\rho,\gamma}(F_x)(c)=0\).  Hence the
  entire weight-\(\gamma\) part cancels and
  \(\Val(F(x'))>\gamma\).
\end{proof}

At a normalized stage of multiplicity \(m\), condition
\[
  \Val(b_i)+i\rho\ge\gamma\qquad(0<i<m)
  \tag{NB}\label{eq:successor-nonbelow}
\]
together with integrality above degree \(m\) puts every weight on or above
\(\gamma\).  The weight-level polynomial has odd degree \(m\) and nonzero
constant term, so \namedref{Lemma}{lem:proper-odd-root} supplies the nonzero
root \(c\) required by \namedref{Lemma}{lem:successor-correction}.

For \(m>1\), a \emph{fixed-lift, fixed-multiplicity Newton chain} is a
natural-number sequence \((x_n)\) in \(\mathfrak m\), using one fixed
integral coefficient lift of \(F\), in which every successor is a correction
from \namedref{Lemma}{lem:successor-correction}.  In
\(x_{n+1}-x_n=c_nt^{\rho_n}\), we call \(\rho_n\) the
\emph{correction value} (or correction exponent).

\begin{lemma}[Strict increase of successive correction values]
\label{lem:successor-values-strict}
  In a fixed-lift, fixed-multiplicity Newton chain with \(m>1\), the map
  \(n\mapsto\rho_n\) is strictly increasing.
\end{lemma}
\begin{proof}
  The successor correction gives
  \(\Val(F(x_{n+1}))>\Val(F(x_n))\).  Writing these values as
  \(m\rho_{n+1}\) and \(m\rho_n\), respectively, and using that
  multiplication by the positive integer \(m\) is strictly order preserving
  gives the claim.
\end{proof}

\begin{lemma}[Hahn sum of a separated family]
\label{lem:ordinal-hahn-sum}
  Let \(I\) be a well-ordered set.  For every \(i\in I\), let \(D_i\) be a
  Hahn series and let \(\ell_i,u_i\in\Gamma\).  Assume
  \[
    \Supp(D_i)\subseteq[\ell_i,u_i),
    \qquad i<j\Longrightarrow u_i\le\ell_j.
    \tag{SS}\label{eq:separated-supports}
  \]
  Then \((D_i)_{i\in I}\) is Hahn summable.
\end{lemma}
\begin{proof}
  To see that \(U=\bigcup_i\Supp(D_i)\) is well ordered, take a nonempty
  subset \(A\subseteq U\).  Choose the least index \(i_0\) for which
  \(A\cap\Supp(D_{i_0})\ne\varnothing\), and then the least element \(a_0\)
  of that intersection.  No earlier support meets \(A\), and
  \eqref{eq:separated-supports} puts every element of every later support at
  least \(u_{i_0}>a_0\).  Thus \(a_0\) is the least element of \(A\).
  Moreover a fixed exponent cannot belong to two distinct supports: if
  \(i<j\), membership in both would give
  \(\gamma<u_i\le\ell_j\le\gamma\).  Hence only finitely many (in fact at
  most one) summands contribute to each coefficient, proving Hahn
  summability.
\end{proof}

In particular, if \(\lambda\) is a limit ordinal and
\((\rho_\alpha)_{\alpha<\lambda}\) is strictly increasing, apply the lemma
with \(D_\alpha=c_\alpha t^{\rho_\alpha}\),
\(\ell_\alpha=\rho_\alpha\), and
\(u_\alpha=\rho_{\alpha+1}\).  Since \(\alpha+1<\lambda\), this proves that
\(x_0+\sum_{\alpha<\lambda}c_\alpha t^{\rho_\alpha}\) is a Hahn series.

For a Hahn-summable family \(A=(A_*)\cup(A_i)_{i\in I}\), a polynomial
\(G(T)=\sum_{j=0}^d b_jT^j\), and \(\gamma\in\Gamma\), define
\(E_{G,\gamma}(A)\) as follows.  Expand each coefficient
\([t^\gamma]b_j(\sum A)^j\); collect every ordinary index \(i\in I\) that
occurs in a tuple with nonzero contribution, and take the union over
\(0\le j\le d\).  Hahn summability makes this a finite set.  The
distinguished index \(*\) is not included in \(E_{G,\gamma}(A)\).

\begin{lemma}[Finite dependence of polynomial-evaluation coefficients]
\label{lem:ordinal-eval-finite}
  Let \(A=(A_*)\cup(A_i)_{i\in I}\) be a Hahn-summable family, where the
  distinguished term \(A_*\) is always retained.  For a predicate
  \(p\) on \(I\), write \(A|_p\) for the family obtained by replacing
  \(A_i\) by zero when \(p(i)\) is false.  Given
  \(G\in k((\Gamma))[T]\) and \(\gamma\in\Gamma\), one has
  \[
    E_{G,\gamma}(A)\subseteq\{i:p(i)\}
    \quad\Longrightarrow\quad
    [t^\gamma]G\!\left(\sum A\right)
      =[t^\gamma]G\!\left(\sum(A|_p)\right).
    \tag{FD}\label{eq:finite-dependence}
  \]
\end{lemma}
\begin{proof}
  By the preceding definition, if \(p\) contains
  \(E_{G,\gamma}(A)\), truncating away indices outside
  \(p\) removes none of the contributing tuples, so every summand in the
  coefficient expansion is unchanged.  This proves
  \eqref{eq:finite-dependence}.  For an ordinal correction family, any
  initial segment containing this finite set is therefore sufficiently late.
\end{proof}

On the multiplicity-\(1\) side, a whole \(\omega\)-block of natural-number
corrections, rather than one monomial correction, is used as one
stage of the transfinite recursion.  Fix \(F\in\Ocal[X]\).  The block data
below make no coefficient or multiplicity assumption on \(F\); those
assumptions enter only when the data are constructed.  For a nonroot
\(x\in\mathfrak m\), meaning \(F(x)\ne0\), write
\(v(x)=\Val(F(x))\in\Gamma\).

A \emph{Newton block successor} is a move from a nonroot
\(x\in\mathfrak m\) to a nonroot \(x^+\in\mathfrak m\) such that
  \[
    v(x)<v(x^+),\qquad
    \Supp(x^+-x)\subseteq [v(x),v(x^+)).
  \]
It is the compressed datum of the limit of one bounded natural correction
chain.  A further correction from that limit is used only to certify that the
limit is a rootless state from which the construction can continue; that
further correction is not part of the block difference.

\begin{lemma}[A root or a multiplicity-\(1\) block successor]
\label{lem:one-block-successor}
  Let \(F\in\Ocal[X]\), assume its constant coefficient lies in
  \(\mathfrak m\) and its linear coefficient is a unit, and let
  \(x\in\mathfrak m\) be a nonroot.  Then either \(F\) has a root in
  \(\mathfrak m\), or there is a Newton block successor from \(x\).
\end{lemma}
\begin{proof}
  Translate by \(x\).  Its constant coefficient is \(F(x)\in\mathfrak m\),
  and its linear coefficient is congruent modulo \(\mathfrak m\) to the
  original unit coefficient; all contributions from higher degrees contain a
  positive-valuation power of \(x\).  Thus the translated polynomial is again
  a normalized multiplicity-\(1\) problem.  Unless a root occurs, repeatedly
  apply \namedref{Lemma}{lem:successor-correction}.  This gives corrections
  \(c_nt^{\rho_n}\) with strictly increasing positive values and finite
  approximations
  \(x_n=x+\sum_{j<n}c_jt^{\rho_j}\), where
  \(\Val(F(x_n))=\rho_n\).

  The correction monomials are Hahn summable; put
  \(x_\omega=x+\sum_nc_nt^{\rho_n}\).  The tail
  \(x_\omega-x_{n+1}\) has valuation \(\rho_{n+1}\).  Taylor expansion, with
  integral coefficients and integral arguments, gives
  \[
    \Val\bigl(F(x_\omega)-F(x_{n+1})\bigr)\ge\rho_{n+1}.
  \]
  Since \(\Val(F(x_{n+1}))=\rho_{n+1}\), it follows that
  \(\Val(F(x_\omega))\ge\rho_{n+1}>\rho_n\).  If the \(\rho_n\) are unbounded, this says
  every coefficient of \(F(x_\omega)\) is zero, so \(x_\omega\) is the first
  alternative.  Otherwise, if no root has occurred, write the finite positive
  value of \(F(x_\omega)\) as \(\eta\).  It satisfies
  \(\rho_n<\eta\) for every \(n\).

  At multiplicity \(1\), the non-below condition is empty.  For the
  translated polynomial at \(x_\omega\), take correction value
  \(\delta=\eta\).  Its weight-level polynomial has nonzero constant and
  linear terms and hence a nonzero root.  One further application of
  \namedref{Lemma}{lem:successor-correction} therefore gives a correction
  from \(x_\omega\).  In the Lean construction this extra step certifies the
  bounded omega-block, but its target is not the block endpoint.

  Set \(x^+=x_\omega\).  It is a nonroot in \(\mathfrak m\),
  \(v(x^+)=\eta>\rho_0=v(x)\), and the support of \(x^+-x\) consists of the
  \(\rho_n\)'s.  Since \(\rho_0\le\rho_n<\eta\),
  \[
    \Supp(x^+-x)\subseteq[\Val(F(x)),\Val(F(x^+))).
  \]
  Thus \(x\mapsto x^+\) is the required Newton block successor.
\end{proof}

\begin{definition}[Coherent multiplicity-\(1\) Newton block process]
\label{def:coherent-one-block}
  Let \(I\) be a well-ordered set with least element \(0\).  A
  \emph{coherent multiplicity-\(1\) Newton block process} on \(I\) consists,
  for every \(\alpha\in I\), of a nonroot \(x_\alpha\in\mathfrak m\) and a
  Newton block successor \(x_\alpha^+\), such that \(\alpha<\beta\) implies
  \[
    v(x_\alpha^+)\le v(x_\beta),
  \]
  and, for every \(\alpha\),
  \[
    x_\alpha-x_0=\sum_{\xi<\alpha}(x_\xi^+-x_\xi).
    \tag{BC}\label{eq:block-coherence}
  \]
  The first condition separates the supports of distinct blocks; the second
  says that every state is the Hahn sum of all preceding block corrections.
\end{definition}

We shall use the following compatibility package at a nonzero limit ordinal
\(\lambda\).  For every \(\tau<\lambda\), suppose a coherent block process
\(C^{<\tau}\) has been constructed on the stages below \(\tau\).  From
\(C^{<\alpha+1}\), select its final state and final block and call
them \(x_\alpha\) and \(B_\alpha\), and write \(x_\alpha^+\) for the endpoint
of \(B_\alpha\).  The family has \emph{natural limit
compatibility} when the following conditions hold.
\begin{enumerate}[label=(\roman*)]
  \item There is one \(x_0\in\mathfrak m\) which is the bottom source of every
    nonempty \(C^{<\tau}\), and the selected state at index \(0\) also has
    source \(x_0\).
  \item If \(\alpha<\beta<\lambda\), the value at the endpoint of
    \(B_\alpha\) equals the endpoint value of the corresponding block after
    \(C^{<\beta+1}\) is restricted through \(\alpha\).
  \item If \(\alpha<\beta<\lambda\), the Hahn sum of the selected blocks
    before \(\alpha\) equals the Hahn sum of the restriction of
    \(C^{<\beta+1}\) before \(\alpha\).
  \item The latter restricted sum equals the prefix sum formed inside
    \(C^{<\alpha+1}\) itself.
\end{enumerate}
Conditions (iii) and (iv) are genuine compatibility assumptions; support
separation alone does not imply that independently constructed prefixes have
the same coefficients.

On the higher-multiplicity side, the transfinite state retains the bounded
natural correction chain itself.  For a normalized problem \(F\) of
multiplicity \(m\), an \emph{integral coefficient lift} is a polynomial
\(D\in\Ocal[X]\) which equals \(F\) in the Hahn field, whose coefficients
below \(m\) lie in \(\mathfrak m\), and whose degree-\(m\) coefficient is a
unit.  With one such \(D\) fixed, a \emph{Newton correction step} from
\(A\in\mathfrak m\) to \(A'\in\mathfrak m\) consists of
\(\gamma,\delta\in\Gamma\) and \(c\in k\) such that
\[
  \gamma>0,\quad\delta>0,\quad c\ne0,\quad m\delta=\gamma,
  \quad A'=A+ct^\delta,
\]
\[
  \Val(D(A))=\gamma,\qquad
  \Val(A'-A)=\delta,\qquad
  \Val(D(A'))>\gamma.
  \tag{NS}\label{eq:newton-step-data}
\]

The \emph{root-or-step property} says the following for every normalized
same-multiplicity zero branch \((F,m)\), after roots at all multiplicities
below \(m\) have been supplied: there is an integral coefficient lift \(D\)
such that, from every \(A\in\mathfrak m\), either \(F\) already has a root in
\(\mathfrak m\), or there is a Newton correction step from \(A\), using that
same \(D\).  Its \emph{fixed-lift form} requires this alternative for every
integral coefficient lift \(D\).  The two forms agree here: two integral
lifts have the same image in the Hahn field, and the inclusion
\(\Ocal\hookrightarrow k((\Gamma))\) is injective, so the lifts are equal.

A sequence \((x_n)_{n\in\mathbb N}\) is a \emph{rootless natural correction
chain} if \(x_0=0\), all \(x_n\) lie in \(\mathfrak m\), and one integral
coefficient lift \(D\) makes every move \(x_n\to x_{n+1}\) a Newton
correction step.  It is \emph{bounded} when its correction values are bounded
above in \(\Gamma\).  A \emph{terminal obstruction} is a same-multiplicity
zero branch equipped with such a bounded chain and its fixed lift \(D\).

A \emph{shifted terminal state} over the original polynomial \(F\) consists
of a shift \(s\in\mathfrak m\), the translated polynomial
\(F_s(Z)=F(s+Z)\), a terminal obstruction for \(F_s\), and the transport
taking a root \(z\) of \(F_s\) to the root \(s+z\) of \(F\).  The first correction value of the
natural chain in this state is called its \emph{rank}.

\begin{definition}[Coherent higher-multiplicity terminal-state process]
\label{def:coherent-terminal-process}
  Let \(I\) be a linearly ordered set, let \(F\in k((\Gamma))[X]\), and let
  \(m\in\mathbb N\).  Write \((S_\alpha)_{\alpha\in I}\) for shifted
  terminal states for \((F,m)\), \(s_\alpha\) for their shifts, and
  \(\rho_\alpha\) for their ranks.  (The terminal data in each state already
  imply that \(m>1\) is odd.)  The family is
  \emph{coherent} when \(\alpha<\beta\) implies
  \(\rho_\alpha<\rho_\beta\), and whenever \(\alpha\le\beta\) and
  \(\gamma<\rho_\alpha\),
  \[
    [t^\gamma]s_\alpha=[t^\gamma]s_\beta.
    \tag{TC}\label{eq:terminal-coherence}
  \]
  Stability of polynomial-evaluation coefficients is not included as a
  field of this definition; it is derived as the recursion invariant below.
\end{definition}

\begin{proposition}[The multiplicity-\(1\) limit prefix]
\label{prop:one-block-limit-prefix}
  Let \(\lambda\) be a nonzero limit ordinal.  Suppose that a process from
  \namedref{Definition}{def:coherent-one-block} has been constructed on every
  proper initial segment and that the resulting family has natural limit
  compatibility.  Then the selected states \(x_\alpha\) and blocks
  \(B_\alpha\), for \(\alpha<\lambda\), form a coherent Newton block process
  on the whole initial segment below \(\lambda\).
\end{proposition}
\begin{proof}
  By (ii), the endpoint value of an earlier selected block is the endpoint
  value of its copy in every later process.  The separation condition in the
  later process therefore gives
  \(v(x_\alpha^+)\le v(x_\beta)\) whenever
  \(\alpha<\beta\).  Hence the selected blocks have separated supports and
  are Hahn summable by \namedref{Lemma}{lem:ordinal-hahn-sum}.

  It remains to check source coherence.  Fix \(\alpha<\lambda\) and choose
  \(\beta\) with \(\alpha<\beta<\lambda\), which is possible because
  \(\lambda\) is a limit ordinal.  Coherence inside \(C^{<\beta+1}\) identifies
  the difference between its state at \(\alpha\) and its bottom source with
  its restricted prefix sum.  By (i) the bottom source is \(x_0\), by (iii)
  the restricted sum is the selected-block sum before \(\alpha\), and by
  (iv) this is also the prefix already computed in \(C^{<\alpha+1}\).  Thus
  \[
    x_\alpha-x_0=\sum_{\xi<\alpha}(x_\xi^+-x_\xi),
  \]
  which is exactly \eqref{eq:block-coherence}.  The state, block-successor,
  and nonroot fields are inherited from the selected final data.  This
  constructs the required coherent process.
\end{proof}

\begin{proposition}[Invariant of the higher-multiplicity terminal process]
\label{prop:terminal-process-invariant}
  Fix a terminal obstruction for \((F,m)\); its branch data imply that
  \(m>1\) is odd.  Assume that every normalized problem of multiplicity
  below \(m\) has been solved, that the fixed-lift root-or-step property
  holds, and that \(F\) has no root in \(\mathfrak m\).  Let
  \((S_\beta)_{\beta\in\mathrm{Ord}}\) be the recursively defined family,
  with shifts \(s_\beta\) and ranks \(\rho_\beta\).  Then for every ordinal
  \(o\):
  \begin{enumerate}[label=(\roman*)]
    \item \(\beta<o\) implies \(\rho_\beta<\rho_o\);
    \item if \(\beta\le o\) and \(\gamma<\rho_\beta\), then
      \([t^\gamma]s_o=[t^\gamma]s_\beta\);
    \item if \(\beta\le o\) and \(\gamma<m\rho_\beta\), then
      \[
        [t^\gamma]F(s_o)=[t^\gamma]F(s_\beta);
        \tag{EC}\label{eq:terminal-eval-coherence}
      \]
    \item the map \(\beta\mapsto\rho_\beta\) on \(\beta<o\) is strictly
      increasing and
      \[
        s_o=\sum_{\xi<o}(s_{\xi+1}-s_\xi).
        \tag{HS}\label{eq:terminal-prefix-sum}
      \]
  \end{enumerate}
  In particular, restriction to any ordinal initial segment is coherent in
  the sense of \namedref{Definition}{def:coherent-terminal-process}.
\end{proposition}
\begin{proof}
  Use transfinite induction on \(\beta\).  At \(0\), take the original
  terminal obstruction with shift \(s_0=0\); all four assertions are
  immediate.

  Suppose \(S_\alpha\) has been constructed.  A lower edge at its first
  correction value would give a root by
  \namedref{Proposition}{prop:lower-edge} and the solved lower-multiplicity
  problems; translating back would contradict the no-root assumption.
  Thus there is no lower edge.  Advance by the first correction of the
  bounded chain and replace the chain by its tail, expressed relative to
  the new source.  The shift increment is a nonzero monomial of
  value \(\rho_\alpha\), the new chain is bounded, and its first correction
  value is strictly larger than \(\rho_\alpha\).  The Newton weight estimate
  behind \eqref{eq:newton-step-data} also says that this move preserves the
  coefficients of \(F(s_\alpha)\) below \(m\rho_\alpha\).  These facts give
  rank growth, shift-coefficient stability, evaluation-coefficient stability,
  and the prefix-sum identity at \(\alpha+1\).

  Now let \(\beta\) be a nonzero limit ordinal.  By induction the increments
  \(s_{\xi+1}-s_\xi\) are nonzero monomials of the strictly increasing values
  \(\rho_\xi\).  They are Hahn summable by
  \namedref{Lemma}{lem:ordinal-hahn-sum}; define \(s_\beta\) to be their sum.
  This proves \eqref{eq:terminal-prefix-sum}, places \(s_\beta\) in
  \(\mathfrak m\), and gives \eqref{eq:terminal-coherence}: below
  \(\rho_\alpha\), every increment from stage \(\alpha\) onward has zero
  coefficient.  For
  \(\alpha<\beta\) and \(\gamma<m\rho_\alpha\), the finite set from
  \namedref{Lemma}{lem:ordinal-eval-finite} is contained in some initial
  segment \(\eta<\beta\).  Taking \(\eta\ge\alpha\) and using the induction
  hypothesis at \(\eta\) proves
  \eqref{eq:terminal-eval-coherence} at the limit.

  The translated polynomial is still a same-multiplicity zero branch.
  Indeed, translation by \(s_\beta\in\mathfrak m\) sends every coefficient
  below degree \(m\) into \(\mathfrak m\): in the Taylor expansion, the old
  coefficients below \(m\) already lie in \(\mathfrak m\), while every term
  coming from degree at least \(m\) contains a positive power of
  \(s_\beta\).  Its degree-\(m\) coefficient is the old unit plus terms in
  \(\mathfrak m\), and hence remains a unit.  Oddness and \(m>1\) are
  unchanged.  Thus the hypotheses needed at the limit state have not been
  lost.

  Apply the fixed-lift root-or-step property at the translated polynomial
  \(F(X+s_\beta)\).  Under the no-root assumption it supplies a natural
  correction chain.  If its values \(\sigma_n\) were unbounded, its monomial
  corrections would have a Hahn sum \(z\).  Given \(\gamma\in\Gamma\), take
  a stage after all indices in the finite-dependence set for the coefficient
  \([t^\gamma]F(s_\beta+z)\), and then take it still later so that
  \(m\sigma_n>\gamma\).  The coefficient equals that at this finite
  approximation, whose evaluation has valuation \(m\sigma_n>\gamma\), so it
  is zero.  Thus every coefficient vanishes and \(F(s_\beta+z)=0\), a
  contradiction.  The new chain is therefore bounded and supplies the
  terminal state at \(\beta\).

  It remains to prove strict rank growth into the limit.  Given
  \(\alpha<\beta\), choose \(\alpha<\eta<\beta\).  The evaluation-coefficient
  stability just proved says that every coefficient of \(F(s_\beta)\) below
  \(m\rho_\eta=\Val(F(s_\eta))\) is zero.  Therefore
  \(m\rho_\eta\le\Val(F(s_\beta))=m\rho_\beta\).  Hence
  \(\rho_\alpha<\rho_\eta\le\rho_\beta\).  This completes all induction
  clauses.
\end{proof}

\begin{definition}[Hartogs ordinal]
\label{def:hartogs}
  For a set \(X\), let \(\#X\) denote its cardinality and let
  \((\#X)^+\) be the successor cardinal, the least cardinal strictly larger
  than \(\#X\).  For a cardinal \(\kappa\), let \(\operatorname{ord}(\kappa)\)
  be the least ordinal of cardinality \(\kappa\).  Define the
  \emph{Hartogs ordinal} by
  \[
    h(X)=\operatorname{ord}\bigl((\#X)^+\bigr).
  \]
  We write \(|h(X)|\) for its underlying well-ordered set of predecessors.
\end{definition}

The underlying set \(|h(X)|\) has cardinality \((\#X)^+\).  An injection
\(|h(X)|\to X\) would therefore imply
\((\#X)^+\le\#X\), contradicting the defining inequality
\(\#X<(\#X)^+\).  Thus this definition has the usual Hartogs property: it
produces a well-ordered set which cannot inject into \(X\).  Equivalently, it
is the least initial ordinal with that property.  This cardinal argument is
the only set-theoretic fact used below.

\begin{lemma}[Hartogs obstruction]
\label{lem:hartogs}
  There is no strictly increasing map from \(|h(\Gamma)|\) into the ordered
  set \(\Gamma\).
\end{lemma}
\begin{proof}
  A strictly increasing map is injective, while the cardinal argument after
  \namedref{Definition}{def:hartogs} says that no injection from
  \(|h(\Gamma)|\) into the underlying set of \(\Gamma\) exists.
\end{proof}

\begin{lemma}[Resolution of a terminal obstruction]
\label{lem:terminal-prefix-resolution}
  Fix a terminal obstruction for \((F,m)\), assume every normalized odd
  problem of multiplicity below \(m\) has been solved, and assume the
  root-or-step property defined above.  Then the obstruction is resolved:
  \(F\) has a root in \(\mathfrak m\).
\end{lemma}
\begin{proof}
  Suppose, for a contradiction, that the original problem has no root in
  \(\mathfrak m\).  A root of any translated problem would transport to a
  root of the original one, so the root alternative is excluded at every
  stage.  The root-or-step property initially selects some integral lift.
  For any prescribed lift, both lifts map to the same Hahn-field polynomial;
  injectivity of \(\Ocal\hookrightarrow k((\Gamma))\) makes them equal.
  Hence the property has the fixed-lift form required by
  \namedref{Proposition}{prop:terminal-process-invariant}.  Applying that
  proposition to the initial terminal obstruction constructs a shifted terminal state
  \(S_\alpha\) for every \(\alpha\in|h(\Gamma)|\), with strictly increasing
  ranks \(\rho_\alpha\).  These ranks define a strictly increasing map
  \[
    |h(\Gamma)|\longrightarrow\Gamma,
  \]
  contradicting \namedref{Lemma}{lem:hartogs}.  The assumption that no root
  exists is therefore false.
\end{proof}

\begin{proposition}[The multiplicity-one case]
\label{prop:multiplicity-one-core}
  A normalized odd Newton problem of multiplicity one has a solution in
  \(\mathfrak m\).
\end{proposition}
\begin{proof}
  Assume that no root exists.  Zero is then a nonroot state.  At every
  successor ordinal, \namedref{Lemma}{lem:one-block-successor} supplies the
  next block.  At a limit ordinal, the already constructed objects are
  restrictions of this one recursion: their bottom source, corresponding
  endpoint values, and prefix Hahn sums therefore satisfy natural limit
  compatibility.  Apply
  \namedref{Proposition}{prop:one-block-limit-prefix} to obtain the coherent
  process on the whole limit initial segment.  Its separated block
  differences are Hahn summable, so define the limit source by
  \[
    x_\lambda=x_0+
      \sum_{\xi<\lambda}(x_\xi^+-x_\xi).
  \]
  Every block difference has positive support, hence
  \(x_\lambda\in\mathfrak m\).  The separation condition also extends to
  this new source.  For \(\alpha<\lambda\), put \(\beta=\alpha+1<\lambda\).
  Coherence gives \(x_\beta=x_\alpha^+\), and every block in the tail
  \(x_\lambda-x_\beta\) has support at least \(v(x_\beta)\).
  Thus \(\Val(x_\lambda-x_\beta)\ge v(x_\beta)\).  Taylor expansion over
  \(\Ocal\) gives
  \[
    \Val(F(x_\lambda)-F(x_\beta))\ge v(x_\beta),
    \qquad
    \Val(F(x_\lambda))\ge v(x_\beta)=v(x_\alpha^+).
  \]
  Under the standing no-root assumption \(x_\lambda\) is a nonroot, so
  this is the required inequality \(v(x_\alpha^+)\le v(x_\lambda)\).
  Applying \namedref{Lemma}{lem:one-block-successor} at
  \(x_\lambda\) supplies the top block for the next successor initial
  segment.  Transfinite recursion therefore continues for as long as no root
  appears.

  The separation inequality in
  \namedref{Definition}{def:coherent-one-block}, together with strict
  increase inside every block, makes the residual values
  \(\alpha\mapsto v(x_\alpha)\) strictly increasing.  If the construction
  continued through all of \(|h(\Gamma)|\), these values would define a strictly
  increasing map from that ordinal into \(\Gamma\), contradicting
  \namedref{Lemma}{lem:hartogs}.  A root must therefore occur, and all corrections
  have positive value, so that root lies in \(\mathfrak m\).
\end{proof}

\begin{proposition}[Termination in the higher-multiplicity branch]
\label{prop:transfinite-termination}
  Let \(m>1\) be odd and assume that every normalized problem of proper odd
  multiplicity below \(m\) has a solution.  Assume also the root-or-step
  property stated above.  Then a normalized problem whose residue root at
  zero has multiplicity \(m\) has a solution in \(\mathfrak m\).
\end{proposition}
\begin{proof}
  Suppose there were no solution.  Repeated application of
  the root-or-step property gives one integral lift and a rootless natural
  correction chain \((x_n)\), with correction values \(\rho_n\): a lower edge would be
  solved by \namedref{Proposition}{prop:lower-edge} and the proper-multiplicity
  hypothesis, while otherwise \namedref{Lemma}{lem:successor-correction}
  supplies the next term.  Suppose
  first that \((\rho_n)\) is unbounded.  Its strictly increasing correction
  monomials are Hahn summable; write their sum as \(x\).  Fix
  \(\gamma\in\Gamma\).  First go beyond every index in the finite set from
  \namedref{Lemma}{lem:ordinal-eval-finite}, and then go still farther, if
  necessary, to choose \(n\) with \(m\rho_n>\gamma\).  The coefficient in
  \(F(x)\) at \(t^\gamma\) is then the same as that of this finite
  approximation.  But
  \[
    \Val(F(x_n))=m\rho_n>\gamma,
  \]
  so that coefficient is zero.  Since \(\gamma\) was arbitrary,
  \(F(x)=0\), contrary to the assumption.

  Thus the correction values are bounded and the chain, together with the
  fixed coefficient data that produces its corrections, is a terminal
  obstruction.  By
  \namedref{Lemma}{lem:terminal-prefix-resolution}, that obstruction itself
  yields a solution in \(\mathfrak m\), a contradiction.
\end{proof}

\section{Odd roots and real closedness of the Hahn field}
\label{sec:hahn-conclusion}

\begin{lemma}[A slope dominated by the top-degree term]
\label{lem:dominating-slope}
  Assume that \(\Gamma\) is nontrivial.  Let
  \(0\ne F(X)=\sum_i a_iX^i\in k((\Gamma))[X]\), put
  \(n=\deg F\), and write \(\gamma_n=\Val(a_n)\).  There is
  \(\delta\in\Gamma\) such that, for every \(i<n\) with \(a_i\ne0\),
  \[
    (n-i)\delta<\Val(a_i)-\gamma_n.
    \tag{DS}\label{eq:dominating-slope}
  \]
\end{lemma}
\begin{proof}
  There are only finitely many indices \(i<n\) with \(a_i\ne0\).  For each
  of them, divisibility of \(\Gamma\) gives a unique threshold \(q_i\) with
  \[
    (n-i)q_i=\Val(a_i)-\gamma_n.
  \]
  Uniqueness follows because an ordered abelian group is torsion-free.
  Choose \(\varepsilon>0\), which exists because \(\Gamma\) is nontrivial.
  If the set of thresholds is nonempty, take
  \(\delta=\min_i q_i-\varepsilon\); if it is empty, take any \(\delta\).
  Then \(\delta<q_i\) is exactly
  \eqref{eq:dominating-slope}.
\end{proof}

\begin{theorem}[Odd-degree roots in the Hahn field]
\label{thm:hahn-odd-roots}
  Every polynomial in \(k((\Gamma))[X]\) of odd degree has a root in
  \(k((\Gamma))\).
\end{theorem}
\begin{proof}
  First prove by strong induction on odd \(m>0\) that every normalized problem
  of multiplicity \(m\) has a root in \(\mathfrak m\).  The case \(m=1\) is
  \namedref{Proposition}{prop:multiplicity-one-core}.  If \(m>1\), the
  induction hypothesis solves every proper odd multiplicity below \(m\).
  Translate at an arbitrary current approximation in \(\mathfrak m\).  If a
  coefficient lies below the current Newton edge,
  \namedref{Proposition}{prop:lower-edge} gives a root.  Otherwise the weight
  inequalities in \eqref{eq:successor-nonbelow},
  \namedref{Lemma}{lem:proper-odd-root}, and
  \namedref{Lemma}{lem:successor-correction} give the next Newton correction.
  Fixing the integral coefficient lift therefore gives the root-or-step
  property required by
  \namedref{Proposition}{prop:transfinite-termination}, which supplies a root.
  This completes the induction.

  Now let \(F\) have odd degree \(n\).  If \(\Gamma\) is trivial, every Hahn
  series has support at the sole group element, so the Hahn field identifies
  with the real closed coefficient field \(k\); the required root then exists
  in \(k\).  Suppose \(\Gamma\) is nontrivial.  Since \(n\) is odd,
  \(F\ne0\) and \(n>0\).  Write \(F=\sum_i a_iX^i\) and
  \(\gamma_n=\Val(a_n)\).  Apply
  \namedref{Lemma}{lem:dominating-slope} to obtain \(\delta\), put
  \(\lambda=\gamma_n+n\delta\), and define
  \(\widetilde F(X)=t^{-\lambda}F(t^\delta X)\).  Its degree-\(n\)
  coefficient has value \(-\lambda+\gamma_n+n\delta=0\), while for
  \(i<n\) with \(a_i\ne0\) the degree-\(i\) coefficient has value
  \[
    -\lambda+\Val(a_i)+i\delta
      =\Val(a_i)-\gamma_n-(n-i)\delta>0.
  \]
  Coefficients above degree \(n\) vanish.  Thus \(\widetilde F\) is a
  normalized problem of odd multiplicity \(n\).
  The induction just proved gives \(z\in\mathfrak m\) with
  \(\widetilde F(z)=0\).  Therefore
  \[
    F(t^\delta z)=t^\lambda\widetilde F(z)=0,
  \]
  so \(t^\delta z\) is a root of the original polynomial.
\end{proof}

\begin{proof}[Proof of \namedref{Theorem}{thm:hahn-real-closed}]
  \namedref{Proposition}{prop:hahn-squares} gives square roots of all nonnegative
  elements, and \namedref{Theorem}{thm:hahn-odd-roots} gives roots of all
  odd-degree polynomials.  \namedref{Proposition}{prop:real-closed-criterion}
  now applies.
\end{proof}

\section{Newton--Puiseux construction with bounded denominators}
\label{sec:puiseux-process}

For \(N\ge1\), put
\[
  \Gamma_N=\frac1N\mathbb Z,\qquad
  \mathcal P_N(k)=
  \{x\in k((\mathbb Q)):\Supp(x)\subseteq\Gamma_N\}.
\]
Then \(\Puiseux=\bigcup_{N\ge1}\mathcal P_N(k)\).  A finite collection of
Puiseux series and polynomial coefficients is contained in one common
\(\mathcal P_N(k)\).

The argument is the classical Newton--Puiseux construction.  We record the
extra denominator control needed for roots in the union of the rings
\(\mathcal P_N(k)\);
see \cite{Walker}.

We call \(\mathcal P_N(k)\) the \emph{fixed denominator level \(N\)}.  For a
polynomial over its valuation ring, the \emph{residue polynomial} is obtained
by reducing every coefficient modulo the maximal ideal.  A residue root
\(a\in k\) is \emph{simple} when the residue polynomial vanishes at \(a\) but
its derivative does not.  In a Newton step of
current multiplicity \(M\), the \emph{full-multiplicity case} is the case in
which the chosen root of the initial polynomial has multiplicity exactly
\(M\); a smaller multiplicity is proper.

For a local ring with maximal ideal \(\mathfrak n\), its
\emph{maximal-ideal topology} is the topology in which congruence modulo
\(\mathfrak n^r\), for \(r=1,2,\ldots\), gives a neighborhood basis of
zero.

\begin{definition}[Power-series model of a fixed denominator level]
\label{def:fixed-level-power-series}
  For \(N\ge1\), put \(q=t^{1/N}\) and
  \[
    \mathcal O_N=\mathcal P_N(k)\cap\Ocal
  \].
  The \emph{power-series identification} of this level is the coefficient
  re-indexing isomorphism
  \[
    \Phi_N:k[[q]]\xrightarrow{\ \sim\ }\mathcal O_N,
    \qquad
    \sum_{i\ge0}a_iq^i\longmapsto\sum_{i\ge0}a_it^{i/N}.
  \]
  Indeed, the image support is well ordered.  Conversely, every exponent in
  an element of \(\mathcal O_N\) is nonnegative and belongs to
  \(\frac1N\mathbb Z\), hence is uniquely \(i/N\) for some
  \(i\in\mathbb N\).  Re-indexing by this unique \(i\) is the inverse map;
  addition and convolution multiplication are preserved because
  \((i+j)/N=i/N+j/N\).
\end{definition}

\begin{lemma}[Completeness of a fixed denominator level]
\label{lem:fixed-level}
  The ring \(\mathcal O_N\) is complete for its maximal-ideal topology.
\end{lemma}
\begin{proof}
  Under the power-series identification from
  \namedref{Definition}{def:fixed-level-power-series}, the maximal ideal of
  \(\mathcal O_N\) corresponds to \((q)\).  Thus
  congruence modulo \((q)^r\) means equality of the coefficients of
  \(1,q,\ldots,q^{r-1}\).  A Cauchy sequence therefore has each coefficient
  eventually constant.  Taking those eventual values produces a series in
  \(k[[q]]\), and the sequence converges to it modulo every \((q)^r\).
  Transporting this limit through \(\Phi_N\) proves the asserted completeness
  of \(\mathcal O_N\).
\end{proof}

\begin{lemma}[Simple roots lift without changing the denominator]
\label{lem:simple-lift}
  Let \(F\in\mathcal O_N[X]\), and suppose that its residue polynomial has
  a simple root \(a\in k\).  Then \(F\) has a root in \(\mathcal O_N\) reducing
  to \(a\).
\end{lemma}
\begin{proof}
  This is the simple-root form of Hensel's lemma for a complete local ring.
  We include its Newton argument.
  Start with a lift \(x_0\) of \(a\).  Then \(F(x_0)\in(q)\), while
  \(F'(x_0)\notin(q)\), so \(F'(x_0)\) is a unit.  Define Newton's iteration
  \[
    x_{n+1}=x_n-\frac{F(x_n)}{F'(x_n)}
  \]
  and put \(e_n=F(x_n)\).  Taylor expansion in the polynomial ring gives
  \[
    F(x_n+h)=F(x_n)+F'(x_n)h+h^2R_n(h)
  \]
  for a polynomial \(R_n\) over \(\mathcal O_N\).  With
  \(h=-e_n/F'(x_n)\), the first two terms cancel, hence
  \[
    v(e_{n+1})\ge 2v(e_n),\qquad
    v(x_{n+1}-x_n)=v(e_n).
  \]
  Moreover \(x_{n+1}\equiv x_n\pmod{(q)}\), so
  \(F'(x_{n+1})\equiv F'(x_n)\not\equiv0\pmod{(q)}\); every derivative used
  above remains a unit.  Since \(v(e_0)\ge1/N\), induction yields
  \(v(e_n)\ge2^n/N\), and therefore
  \((x_n)\) is Cauchy for the \((q)\)-adic topology.  By
  \namedref{Lemma}{lem:fixed-level} it converges to some \(x\in\mathcal O_N\).
  The inequalities \(v(e_n)\to\infty\) and continuity of polynomial
  evaluation give \(F(x)=0\), while \(x\equiv a\pmod{(q)}\).
\end{proof}

\begin{lemma}[Full multiplicity preserves the denominator lattice]
\label{lem:full-multiplicity-denominator}
  Let \(F\in\Ocal[X]\) and \(y\in\Ocal\), and suppose that the supports of
  every coefficient of \(F\) and of \(y\) lie in \(\Gamma_N\).  Let
  \(M>0\), \(\delta>0\), and \(\gamma=M\delta\).  Assume that the coefficient
  of \(X^M\) in \(F(X+y)\) is a unit and that
  \(\Init_{\delta,\gamma}(F(X+y))\) has a nonzero root
  \(c\in k^\times\) of multiplicity exactly \(M\).  Then
  \(\delta\in\Gamma_N\).
\end{lemma}
\begin{proof}
  For \(i>M\), the translated coefficient is integral and hence its Newton
  weight is at least \(i\delta>M\delta=\gamma\).  Thus the initial polynomial
  has degree at most \(M\), while its degree-\(M\) coefficient is nonzero
  because that translated coefficient is a unit.  Its degree is therefore
  exactly \(M\).  A root of multiplicity \(M\) in a degree-\(M\) polynomial
  forces the factorization
  \[
    I(T)=u(T-c)^M,\qquad u\ne0.
  \]
  Its coefficient of \(T^{M-1}\) is \(-uMc\ne0\).  Since the translated
  degree-\(M\) coefficient is a unit, its exponent on the weight level
  \(\gamma\) is \(0=\gamma-M\delta\).  The coefficient of \(T^{M-1}\)
  comes from exponent
  \(\gamma-(M-1)\delta\).  This exponent is therefore equal to
  \(\delta\).  It belongs to the support of the degree-\((M-1)\)
  coefficient of \(F(X+y)\).  Translation by \(y\) preserves
  \(\Gamma_N\)-supported coefficients, because only finite sums and products
  occur.  Hence this translated coefficient, and therefore \(\delta\), lies
  in \(\Gamma_N\).
\end{proof}

\begin{lemma}[A strict chain in a rational lattice is unbounded]
\label{lem:lattice-chain}
  Let \(f:\mathbb N\to\mathbb Q\).  If \(f\) is strictly increasing and
  \(f(n)\in\Gamma_N\) for every \(n\), then \(f\) is not bounded above.
\end{lemma}
\begin{proof}
  More precisely,
  \[
    f(n+1)-f(n)\ge\frac1N,\qquad
    f(n)\ge f(0)+\frac nN.
  \]
  The difference of two elements of \(\Gamma_N\) is again in \(\Gamma_N\),
  and every positive element of \(\Gamma_N\) is at least \(1/N\).
  Summing gives the second inequality.  If \(b\) were an upper bound, choose
  \(n>N(b-f(0))\); then \(f(0)+n/N>b\), a contradiction.
\end{proof}

\begin{theorem}[Bounded-denominator odd roots]
\label{thm:puiseux-odd-roots}
  Every odd-degree polynomial over \(\Puiseux\) has a root in \(\Puiseux\).
\end{theorem}
\begin{proof}
  Let \(P\) be the given polynomial.  Put all its coefficients in one
  denominator level.  Write \(d=\deg P\) and let \(\gamma_d\) be the value of
  its leading coefficient.  Apply \namedref{Lemma}{lem:dominating-slope} with
  \(\Gamma=\mathbb Q\), choose \(\delta\), and put
  \(\lambda=\gamma_d+d\delta\).  Enlarge the denominator once so that
  \(\delta,\lambda\in\Gamma_N\).  The same coefficient calculation as in
  the proof of \namedref{Theorem}{thm:hahn-odd-roots} shows that
  \[
    F(X)=t^{-\lambda}P(t^\delta X)
  \]
  is a normalized problem all of whose coefficients lie in
  \(\mathcal P_N(k)\).  We solve \(F\) by strong induction on its odd
  multiplicity.
  A simple residue root is lifted by \namedref{Lemma}{lem:simple-lift}.  A proper
  odd multiplicity \(r<M\) is normalized by one translation and one monomial
  dilation.  Only finitely many new rational monomial exponents occur.
  The scalar constants lie in \(k\) and have support at zero, while the
  existing coefficients already have a common denominator.  Thus, after
  replacing \(N\) by one common multiple of the old and new denominators,
  all transformed coefficients lie in a single denominator level.  The induction hypothesis at \(r\) gives a root
  at some bounded level, and undoing these two finite transformations still
  has bounded denominator.  Thus denominator enlargement in every
  multiplicity-decreasing branch is finite; an infinite branch can only be
  the full-multiplicity branch treated next.

  In the full-multiplicity case of current multiplicity \(M\),
  \namedref{Lemma}{lem:full-multiplicity-denominator}
  shows inductively that the translated polynomial and every successive
  correction remain at the same fixed level \(\Gamma_N\).  If the corrections
  terminate, the finite sum of corrections is a root in
  \(\mathcal P_N(k)\).  Otherwise write the corrections as
  \(c_nt^{\delta_n}\).  The values \(\delta_n\) are strictly increasing in
  \(\Gamma_N\), hence are unbounded by \namedref{Lemma}{lem:lattice-chain}.

  The family \((c_nt^{\delta_n})_n\) is Hahn summable and its sum \(x\)
  still has support in \(\Gamma_N\).  Fix an exponent \(\gamma\).  First go
  beyond every index in the finite set supplied by
  \namedref{Lemma}{lem:ordinal-eval-finite}, and then increase \(n\), if
  necessary, until \(\delta_n>\gamma\).  The coefficient of \(t^\gamma\) in
  \(F(x)\) is then the coefficient at the \(n\)-th approximation.  At that stage the
  successor construction gives \(\Val(F(x_n))=M\delta_n\).  Since \(M>0\)
  and \(\delta_n>\gamma\), this valuation is greater than \(\gamma\), so the
  coefficient of \(t^\gamma\) in \(F(x_n)\), and hence in \(F(x)\), is
  zero.  Thus every coefficient of \(F(x)\) is zero and
  \(F(x)=0\).  Consequently \(x\in\mathcal P_N(k)\), and
  \(t^\delta x\in\Puiseux\) is a root of the original polynomial \(P\).
\end{proof}

\begin{proposition}[Nonnegative Puiseux series are squares]
\label{prop:puiseux-squares}
  Every nonnegative element of \(\Puiseux\) is a square in \(\Puiseux\).
\end{proposition}
\begin{proof}
  The case \(x=0\) is immediate.  If \(0<x\in\mathcal P_N(k)\), the leading-term argument of
  \namedref{Section}{sec:square-roots} halves the leading exponent and the
  binomial series uses only sums of exponents in \(\Gamma_N\).  Thus the
  constructed square root has support in
  \(\Gamma_{2N}\), and hence belongs to \(\Puiseux\).
\end{proof}

\begin{proof}[Proof of \namedref{Theorem}{thm:puiseux-real-closed}]
  Apply \namedref{Proposition}{prop:puiseux-squares},
  \namedref{Theorem}{thm:puiseux-odd-roots}, and
  \namedref{Proposition}{prop:real-closed-criterion}.
\end{proof}

\section{Corollaries with real coefficients}
\label{sec:real}

We use the classical fact that \(\mathbb R\) is real closed.  Equivalently,
every nonnegative real number has a real square root and every odd-degree
real polynomial has a real zero.

\begin{corollary}[The real rational Hahn field]
\label{cor:real-rational-hahn}
  The ordered Hahn field \(\mathbb R((\mathbb Q))\) is real closed.
\end{corollary}
\begin{proof}
  The ordered additive group \(\mathbb Q\) is divisible: for
  \(q\in\mathbb Q\) and \(n>0\), the element \(q/n\) satisfies
  \(n(q/n)=q\).  \namedref{Theorem}{thm:hahn-real-closed}, with
  \(k=\mathbb R\) and \(\Gamma=\mathbb Q\), therefore applies.
\end{proof}

\begin{corollary}[Real Puiseux field]
\label{cor:real-puiseux}
  The field \(\mathcal P(\mathbb R)\) is real closed.
\end{corollary}
\begin{proof}
  Apply \namedref{Theorem}{thm:puiseux-real-closed} with \(k=\mathbb R\).
\end{proof}

\appendix

\section{Correspondence with the Lean formalization}
\label{sec:lean}

The main text is deliberately mathematical and uses no Lean terminology.
This appendix records the correspondence needed to audit that its case splits,
induction parameter, limit construction, and denominator control agree with
the maintained formal proof.  Declaration and module names here are audit
markers, not substitutes for arguments in the main text.  Each assertion in
the left column links back to its complete statement and proof.
\lean{HahnField.exists_leadingTerm_decomposition_of_pos},
\lean{KSameMultiplicityOrdinalJetChain},
\lean{iterateFixedLiftTerminalOrdinalPrefixState_invariant}, and
\lean{not_bddAbove_range_of_strictMono_of_common_denominator} are private
Lean declarations: they are source-level audit markers, not importable global
declaration names.
\begin{longtable}{@{}>{\raggedright\arraybackslash}p{0.52\textwidth}
  >{\raggedright\arraybackslash}p{0.42\textwidth}@{}}
\toprule
\textbf{Mathematical assertion} & \textbf{Formal correspondence and role}\\
\midrule
\endhead
\claimref{thm:hahn-real-closed}{Theorem}{real closedness of \(k((\Gamma))\)}
  & \lean{hahnKaplansky_realClosed_of_realClosed_divisible}.\\
\claimref{thm:puiseux-real-closed}{Theorem}{real closedness of \(\mathcal P(k)\)}
  & \lean{puiseuxSeries_isRealClosed}.\\
\claimref{prop:real-closed-criterion}{Proposition}{the ordered-field criterion for real closedness}
  & \lean{IsRealClosed.of_linearOrderedField}.\\
\claimref{lem:leading-term}{Lemma}{leading-term normalization of a positive Hahn series}
  & \lean{HahnField.exists_leadingTerm_decomposition_of_pos}.\\
\claimref{lem:positive-unit-square}{Lemma}{binomial square root of a \(1\)-unit}
  & \lean{HahnField.exists_square_root_of_orderTop_sub_one_pos}.\\
\claimref{prop:hahn-squares}{Proposition}{nonnegative Hahn series are squares}
  & \lean{HahnField.isSquare_of_nonneg}.\\
\claimref{lem:newton-correction}{Lemma}{normalization at an odd-multiplicity Newton root}
  & \lean{HahnField.oddClusterHypotheses_scalePolynomial_comp_single_zero_of_rootMultiplicity}.\\
\claimref{def:odd-newton-problem}{Definition}{normalized odd Newton problem}
  & \lean{HahnField.KOddClusterHypotheses}.\\
\claimref{lem:proper-odd-root}{Lemma}{a nonzero odd-multiplicity root of an odd-degree polynomial}
  & \lean{exists_nonzero_root_odd_rootMultiplicity_le_natDegree}.\\
\claimref{prop:lower-edge}{Proposition}{multiplicity descent along a lower Newton edge}
  & \lean{HahnField.positiveHahnRoot_of_lowerNewtonEdge}.\\
\claimref{def:ordinal-terminology}{Definition}{ordinal initial segments, successors, and limits}
  & Standard mathlib concepts.\\
\claimref{lem:successor-correction}{Lemma}{successor-stage correction}
  & \lean{HahnField.exists_oddClusterStep_of_current_edge_initial_root}.\\
\claimref{lem:successor-values-strict}{Lemma}{strict increase of successive correction values}
  & \lean{HahnField.KOddClusterFixedStepChainFromData.stepValue_strictMono}.\\
\claimref{lem:ordinal-hahn-sum}{Lemma}{Hahn sum of a separated well-ordered family}
  & \lean{HahnField.HahnSeries.summableFamilyOfSeparatedSupports}.\\
\claimref{lem:ordinal-eval-finite}{Lemma}{finite dependence of polynomial-evaluation coefficients}
  & \lean{HahnField.OrdinalSupport.HahnSummableFamily.eval_coeff_eq_predTrunc_of_support}.\\
\claimref{lem:one-block-successor}{Lemma}{a root or a multiplicity-\(1\) block successor}
  & \lean{HahnField.rootInMaximalIdeal_or_exists_blockSuccessor_at_state}.\\
\claimref{def:coherent-one-block}{Definition}{coherent multiplicity-\(1\) Newton block process}
  & \lean{HahnField.OneClusterCoherentBlockChain}.\\
\claimref{def:coherent-terminal-process}{Definition}{coherent higher-multiplicity terminal-state process}
  & \lean{KSameMultiplicityOrdinalJetChain}.\\
\claimref{prop:one-block-limit-prefix}{Proposition}{the multiplicity-\(1\) limit prefix}
  & \lean{HahnField.PositiveIioCoherentChainAtBelow.limitStepOfNaturalComponents}.\\
\claimref{prop:terminal-process-invariant}{Proposition}{invariant of the higher-multiplicity terminal process}
  & \lean{iterateFixedLiftTerminalOrdinalPrefixState_invariant}.\\
\claimref{def:hartogs}{Definition}{the Hartogs ordinal}
  & \lean{hartogsOrdinal}.\\
\claimref{lem:hartogs}{Lemma}{the Hartogs obstruction in \(\Gamma\)}
  & \lean{not_exists_strictMono_cardinal_succ_ord_toType}.\\
\claimref{lem:terminal-prefix-resolution}{Lemma}{resolution of a terminal obstruction}
  & \lean{HahnField.KSameMultiplicityFixedStepTerminalObstructionData.positiveHahnRoot_of_edgeStep_rootOrStep_terminalResolution}.\\
\claimref{prop:multiplicity-one-core}{Proposition}{solution of the multiplicity-\(1\) branch}
  & \lean{kOddClusterRootAt_one_of_ordinalStackRec}.\\
\claimref{prop:transfinite-termination}{Proposition}{termination of the higher-multiplicity branch}
  & \lean{positiveHahnRoot_of_sameMultiplicityZeroBranch}.\\
\claimref{lem:dominating-slope}{Lemma}{a slope dominated by the top-degree term}
  & \lean{HahnField.exists_natDegree_dominating_slope}.\\
\claimref{thm:hahn-odd-roots}{Theorem}{roots of odd-degree polynomials in \(k((\Gamma))\)}
  & \lean{hahnKaplansky_exists_root_of_odd_natDegree}.\\
\claimref{def:fixed-level-power-series}{Definition}{power-series identification of a fixed denominator level}
  & \lean{Puiseux.powerSeriesFixedDenominatorValuationSubringEquiv}.\\
\claimref{lem:fixed-level}{Lemma}{completeness of a fixed denominator level}
  & \lean{Puiseux.fixedDenominatorValuationSubring_isAdicComplete}.\\
\claimref{lem:simple-lift}{Lemma}{simple-root lifting at a fixed denominator}
  & \lean{Puiseux.levelSimpleRootLiftHypothesis_of_isAdicComplete}.\\
\claimref{lem:full-multiplicity-denominator}{Lemma}{denominator preservation at full multiplicity}
  & \lean{Puiseux.hasDenominator_newtonStep_of_full_initial_rootMultiplicity}.\\
\claimref{lem:lattice-chain}{Lemma}{unboundedness of a strict chain in \(\frac1N\mathbb Z\)}
  & \lean{not_bddAbove_range_of_strictMono_of_common_denominator}.\\
\claimref{thm:puiseux-odd-roots}{Theorem}{bounded-denominator roots of odd-degree polynomials}
  & \lean{puiseuxSeries_exists_root_of_odd_natDegree}.\\
\claimref{prop:puiseux-squares}{Proposition}{nonnegative Puiseux series are squares}
  & \lean{puiseuxSeries_isSquare_of_nonneg}.\\
\claimref{cor:real-rational-hahn}{Corollary}{real closedness of \(\mathbb R((\mathbb Q))\)}
  & \lean{realRationalHahnField_isRealClosed}.\\
\claimref{cor:real-puiseux}{Corollary}{real closedness of \(\mathcal P(\mathbb R)\)}
  & \lean{realPuiseuxSeries_isRealClosed}.\\
\bottomrule
\end{longtable}

\section{Reproducibility and references}
\label{sec:reproducibility}

The machine-checked developments can be rebuilt with:
\begin{verbatim}
lake build
lake build HahnKaplanskyRealClosedness.Examples
lake env lean HahnKaplanskyRealClosedness/AxiomAudit.lean
lake env lean HahnKaplanskyRealClosedness/Smoke.lean
lake env lean HahnKaplanskyRealClosedness/PuiseuxSmoke.lean
task lint
\end{verbatim}
The public theorem audit uses only the standard foundational axioms:
\[
  \texttt{propext},\qquad
  \texttt{Classical.choice},\qquad
  \texttt{Quot.sound}.
\]

\begin{thebibliography}{9}
\bibitem{Kaplansky}
  I. Kaplansky, \emph{Maximal fields with valuations},
  Duke Mathematical Journal 9 (1942), 303--321.
\bibitem{EnglerPrestel}
  A. J. Engler and A. Prestel, \emph{Valued Fields},
  Springer Monographs in Mathematics, Springer, 2005.
\bibitem{Walker}
  R. J. Walker, \emph{Algebraic Curves},
  Princeton Mathematical Series 13, Princeton University Press, 1950.
\end{thebibliography}

\end{document}
